% Laget automatisk av elevark/lag_elevark.py fra opplegg/energi-gravitasjonsfelt.html. Endringer her blir overskrevet.
\documentclass[11pt]{article}
\usepackage{elevark}
\begin{document}
\elevarkhode{Fysikk 2 · Gravitasjon}{Energi i gravitasjonsfeltet – arbeidsark}{https://fysikklab.no/sim/unnslipningsfart.html}{Unnslipningsfart}

\section*{Oppgaver}

Oppgave 1–4 bruker Unnslipningsfart\footnote{\url{fysikklab.no/sim/unnslipningsfart.html}}, og oppgave 5–7 bruker Satellittbaner\footnote{\url{fysikklab.no/sim/satellittbaner.html}}. Oppgave 8–10 er uten simulering. Bruk $γ = 6{,}674 \cdot 10^{-11}\,\mathrm{N\,m^{2}/kg^{2}}$, jordas masse $5{,}97 \cdot 10^{24}\,\mathrm{kg}$ og jordas radius $6{,}37 \cdot 10^{6}\,\mathrm{m}$.

\begin{oppgaveliste}
\oppgave \fase{gjett}{Gjett}En rakett skytes rett opp fra jorda med 5,0~km/s. Regn ut høyden med mgh. Kommer raketten høyere eller lavere i virkeligheten? Sjekk med simuleringen.
\svar
\oppgave \fase{test}{Regn}Bruk energibevaring med $E_{\mathrm{p}} = -γMm/r$ og regn ut hvor høyt raketten kommer med 5,0~km/s og med 8,0~km/s.
\svar
\oppgave \fase{forklar}{Forklar}Se på grafen i simuleringen. Hva betyr det at linja for samlet energi ligger under null? Hva skjer når den ligger over null?
\svar
\oppgave \fase{test}{Regn}Vis at unnslipningsfarten er $\sqrt{2γM/R}$, og regn den ut for jorda og månen ($7{,}35 \cdot 10^{22}\,\mathrm{kg}$, $1{,}74 \cdot 10^{6}\,\mathrm{m}$). Sjekk med simuleringen.
\svar
\oppgave \fase{test}{Regn}Vis at en satellitt i sirkelbane har $E_{\mathrm{k}} = -\tfrac{1}{2}E_{\mathrm{p}}$. Sjekk med energigrafen for romstasjonen i 400~km høyde.
\svar
\oppgave \fase{test}{Undersøk}Gi satellitten 110~\% av sirkelbanefarten. Hva skjer med $E_{\mathrm{k}}$ og $E_{\mathrm{p}}$ rundt banen? Hvorfor er den samlede energien konstant?
\svar
\oppgave \fase{test}{Undersøk}Finn hvor mange prosent av sirkelbanefarten som trengs for at satellitten skal slippe unna. Forklar svaret med formlene.
\svar
\oppgave \fase{forklar}{Forklar}Vis at $ΔE_{\mathrm{p}} = γMm(1/R - 1/(R + h))$ blir omtrent mgh når $h$ er mye mindre enn $R$.
\svar
\oppgave \fase{test}{Regn}Hvor stort arbeid må gjøres for å løfte 1,0~kg fra bakken til 400~km høyde? Sammenlign med mgh. Hvor mange prosent feil gir mgh?
\svar
\oppgave \fase{utfordring}{Utfordring}En satellitt på 1000~kg skal fra bakken til en sirkelbane i 400~km høyde. Hvor mye energi trengs til sammen? Hvor stor del går til å løfte den, og hvor stor del til å gi den fart? Se bort fra jordrotasjonen og luftmotstanden.
\svar
\end{oppgaveliste}

\end{document}
